\item \textbf{Problema de repaso.}
Para el arreglo que se muestra en la figura P17.40, el plano inclinado y la pequeña polea no tienen fricción, la cuerda soporta el objeto de masa $M$ en el fondo del plano y la cuerda tiene masa $m$. El sistema está en equilibrio y la parte vertical de la cuerda tiene una longitud $h$. Queremos estudiar las ondas que se establecen en la sección vertical de la cuerda.
\begin{enumerate}
    \item ¿Qué análisis de modelo describe al objeto de masa $M$?
    \item ¿Qué análisis de modelo describe las ondas en la parte vertical de la cuerda?
    \item Encuentre la tensión en la cuerda.
    \item Modele la forma de la cuerda como un cateto y la hipotenusa de un triángulo rectángulo. Encuentre la longitud total de la cuerda.
    \item Determine la masa por unidad de longitud de la cuerda.
    \item Encuentre la rapidez de las ondas en la cuerda.
    \item Halle la frecuencia más baja para una onda estacionaria.
    \item Evalúe este resultado para $M = 1.50\text{ kg}$, $m = 0.750\text{ g}$, $h = 0.500\text{ m}$ y $\theta = 30.0^\circ$.
    \item Determine el valor numérico de la frecuencia más baja para una onda en la sección inclinada de la cuerda.
\end{enumerate}